% =============================================================================
% APPENDIX III: METHOD-LLMMC: LLM Monte Carlo Calibration Pipeline
% =============================================================================
% Module: BCM2_14_LLMMC (LLM-based Effect Estimation)
% Category: METHOD
% Version: 1.0 (January 2026)
% Status: Final
% Language: English
%
% This appendix documents the complete LLMMC (LLM Monte Carlo) calibration
% pipeline for converting LLM effect estimates into calibrated α-parameters
% and R-Scores for intervention prioritization decisions.
%
% =============================================================================

% -----------------------------------------------------------------------------
% HEADER BLOCK (Required)
% -----------------------------------------------------------------------------

\begin{tcolorbox}[colback=gray!5!white, colframe=gray!75!black,
    title=Appendix Header: III METHOD-LLMMC]

\begin{tabular}{@{}ll@{}}
\textbf{Code:} & III \\
\textbf{Category:} & METHOD (Methodology) \\
\textbf{Full Name:} & METHOD-LLMMC: LLM Monte Carlo Calibration Pipeline \\
\textbf{Module:} & BCM2\_14\_LLMMC (LLM-based Effect Estimation) \\
\textbf{Version:} & 1.0 (January 2026) \\
\textbf{Status:} & Final \\
\textbf{Language:} & English \\
\end{tabular}

\vspace{0.5em}
\textbf{Purpose:} Complete end-to-end pipeline for converting LLM effect estimates
into calibrated $\alpha$-parameters and R-Scores. Enables GO/PILOT/NO-GO decisions
for intervention design based on empirically anchored calibration.

\end{tcolorbox}

% =============================================================================
% CROSS-REFERENCE STRUCTURE
% =============================================================================
%
% CHAPTER LINKAGE (Required):
% - Primary Chapter: Chapter 17: Intervention Foundations
% - Secondary Chapters: Ch. 20 (Portfolios), Ch. 14 (Segments)
%
% DEPENDENCIES (what this appendix REQUIRES):
% - Appendix AN (METHOD-LLMMC): Core LLMMC theory (LLMs as stochastic estimators)
% - Appendix HHH (METHOD-TOOLKIT): Intervention design schema
% - Appendix EST (CORE-WHERE): Parameter repository for anchors
% - Appendix CTW (CORE-WHEN): Context framework for Ψ
% - Appendix DSN (METHOD-DESIGN): Model design workflow
%
% DEPENDENTS (what REQUIRES this appendix):
% - Appendix HHH: Uses R-Score for intervention prioritization
% - Appendix DSN: Uses calibrated α for model parameters
% - Chapter 17-20: Intervention design block
%
% NOTE: SW (METHOD-SWSM) now references AN directly for LLMMC theory
%
% =============================================================================

% -----------------------------------------------------------------------------
% CHAPTER-APPENDIX LINKAGE BOX
% -----------------------------------------------------------------------------

\begin{tcolorbox}[colback=purple!5!white, colframe=purple!75!black,
    title=Chapter Linkage]

\textbf{Primary Chapter:} Chapter 17: Intervention Foundations
\begin{itemize}[nosep]
\item Section~\ref{sec:17-effects}: Effect estimation $\leftrightarrow$ This Appendix Section~\ref{sec:iii-calibration}
\item Section~\ref{sec:17-prioritization}: Intervention prioritization $\leftrightarrow$ This Appendix Section~\ref{sec:iii-rscore}
\end{itemize}

\textbf{Secondary Chapters:}
\begin{itemize}[nosep]
\item Chapter 20: Intervention Portfolios --- R-Score aggregation
\item Chapter 14: Behavioral Segments --- Context-specific calibration
\end{itemize}

\textbf{Reading Guide:}
\begin{itemize}[nosep]
\item For conceptual overview of $\alpha$: Read \texttt{docs/methods/alpha-definition.md}
\item For $\alpha$ vs. $U$ distinction: Read \texttt{docs/methods/alpha-vs-utility-formal.md}
\item For worked example: Read \texttt{docs/examples/green-energy-default-proof-of-concept.md}
\item For implementation: See \texttt{scripts/llmmc\_to\_rscore.py}
\end{itemize}

\end{tcolorbox}

% -----------------------------------------------------------------------------
% CHAPTER TITLE + LABEL
% -----------------------------------------------------------------------------

\chapter{METHOD-LLMMC: LLM Monte Carlo Calibration Pipeline}
\label{app:iii-llmmc}

% -----------------------------------------------------------------------------
% ABSTRACT
% -----------------------------------------------------------------------------

\begin{abstract}
How can we convert LLM-generated effect estimates into reliable, calibrated parameters
for intervention design? This appendix presents the LLMMC (LLM Monte Carlo) pipeline:
a complete methodology for (1) eliciting effect estimates from LLMs, (2) calibrating
them against Tier-1/2 empirical anchors, (3) mapping to a unified $\theta \in [0,1]$ scale,
(4) computing R-Scores with Monte Carlo uncertainty propagation, and (5) generating
GO/PILOT/NO-GO decisions. The pipeline is empirically validated with 5/5 deployment
checks passed for both d/SMD and percentage-point scales.
\end{abstract}

\textbf{Context:} This appendix covers LLMMC (Tier 3 calibration). For the complete parameterization methodology including all three tiers and method selection, see \textbf{Appendix PM (REF-PARAMETH)}.

% -----------------------------------------------------------------------------
% QUICK REFERENCE BOX
% -----------------------------------------------------------------------------

\begin{tcolorbox}[colback=blue!5!white, colframe=blue!75!black,
    title=Quick Reference: LLMMC Pipeline]

\textbf{Core Question:} How do we convert LLM estimates into calibrated intervention scores?

\textbf{Key Formulas:}

\textbf{Calibration (scale-specific):}
\begin{align}
\text{d/SMD:} \quad \theta_{\text{true}} &= 0.03 + 0.79 \times \theta_{\text{llm}} \label{eq:cal-d} \\
\text{pp:} \quad \text{pp}_{\text{true}} &= -0.48 + 1.007 \times \text{pp}_{\text{llm}} \label{eq:cal-pp}
\end{align}

\textbf{R-Score:}
\begin{equation}
R(a,K) = \sum_{i=1}^{6} \alpha_i + \gamma \cdot \|a\| \cdot \|K\|
\label{eq:rscore}
\end{equation}

\textbf{Decision Rule:}
\begin{equation}
\text{Decision} = \begin{cases}
\text{NO-GO} & \text{if } P(R > T) < 5\% \\
\text{PILOT} & \text{if } 5\% \leq P(R > T) < 25\% \\
\text{GO} & \text{if } P(R > T) \geq 25\%
\end{cases}
\label{eq:decision}
\end{equation}

\textbf{Key Concepts:}
\begin{itemize}[nosep]
\item \textbf{$\alpha_i$}: Calibrated fit parameter for dimension $i$ (not an effect!)
\item \textbf{$\gamma$}: Complementarity between intervention and context
\item \textbf{R-Score}: Relevance/prioritization score (not causal effect)
\item \textbf{EU}: Elicitation Uncertainty (propagated via Monte Carlo)
\end{itemize}

\textbf{Implementation:} \texttt{scripts/llmmc\_to\_rscore.py}

\end{tcolorbox}

% -----------------------------------------------------------------------------
% APPENDIX SCOPE BOX
% -----------------------------------------------------------------------------

\begin{tcolorbox}[
    colback=orange!5!white,
    colframe=orange!75!black,
    title={\textbf{Appendix Scope: Ziel / In-Scope / Out-of-Scope / Constraints}},
    fonttitle=\bfseries
]

\textbf{Ziel (Objective):}
\begin{itemize}[nosep]
    \item How do we know if an intervention will work in a specific context,
before running an expensive RCT?
\end{itemize}

\textbf{In-Scope:}
\begin{itemize}[nosep]
    \item Motivation: The Fundamental Question
    \item The Six $\alpha$-Dimensions
    \item Core Theory: The LLMMC Approach
    \item Calibration Methodology
    \item Piecewise Linear Mapping
\end{itemize}

\textbf{Out-of-Scope (delegiert an andere Appendices):}
\begin{itemize}[nosep]
    \item METHOD-TOOLKIT $\rightarrow$ Appendix HHH
    \item CORE-WHERE $\rightarrow$ Appendix EST
    \item CORE-WHEN $\rightarrow$ Appendix CTW
    \item METHOD-DESIGN $\rightarrow$ Appendix DSN
    \item Central Glossary $\rightarrow$ Appendix GLS
\end{itemize}

\textbf{Constraints (Anwendungsgrenzen):}
\begin{itemize}[nosep]
    \item [TODO: Define when this appendix does NOT apply]
    \item [TODO: Define required prerequisites]
    \item [TODO: Define known limitations]
\end{itemize}

\textbf{Lieferobjekte (Deliverables):}
\begin{itemize}[nosep]
    \item 7 Table(s)
    \item 6 Equation(s)
    \item 7 Axiom(s)
    \item Worked Example(s)
\end{itemize}

\end{tcolorbox}


% =============================================================================
% PART 2: CORE CONTENT
% =============================================================================

\section{Motivation: The Fundamental Question}
\label{sec:iii-motivation}

\begin{tcolorbox}[colback=yellow!5!white, colframe=orange!75!black,
    title=The Central Question]
\textbf{How do we know if an intervention will work in a specific context,
before running an expensive RCT?}
\end{tcolorbox}

The challenge: Behavioral interventions show highly heterogeneous effects across contexts.
A default that produces 69pp effect in one setting may produce 5pp in another.
Traditional meta-analyses provide average effects but not context-specific predictions.

\textbf{Our solution:} Use LLMs to synthesize contextual information from the literature,
then calibrate these estimates against empirical anchors to produce reliable predictions.

\subsection{What $\alpha$ Is (and What It Is Not)}
\label{sec:iii-alpha-definition}

\begin{tcolorbox}[colback=green!5!white, colframe=green!75!black,
    title=The $\alpha$-Parameter: A Fit Measure]

$\alpha$ answers the question:
\begin{quote}
\emph{``If I deploy this intervention here and now --- does it fit?''}
\end{quote}

\textbf{$\alpha$ IS:}
\begin{itemize}[nosep]
\item A calibrated \textbf{fit parameter} (intervention $\times$ context)
\item Empirically anchored (via Tier-1/2 studies)
\item Context-specific (changes with $\Psi$)
\end{itemize}

\textbf{$\alpha$ is NOT:}
\begin{itemize}[nosep]
\item An effect size (effects come from RCTs)
\item A utility measure (utility is normative)
\item A preference (preferences are subjective)
\end{itemize}

\end{tcolorbox}

\textbf{Key insight:}
\begin{quote}
\emph{``Effects say that something works. $\alpha$ says where and when.''}
\end{quote}

% -----------------------------------------------------------------------------
% SECTION: THE 6 ALPHA DIMENSIONS
% -----------------------------------------------------------------------------

\section{The Six $\alpha$-Dimensions}
\label{sec:iii-alpha-dimensions}

Each intervention is evaluated on six fit dimensions:

\begin{table}[h]
\centering
\begin{tabular}{clll}
\toprule
\textbf{Dim} & \textbf{Name} & \textbf{Question} & \textbf{Scale} \\
\midrule
$\alpha_1$ & Phase-Fit & Where in the journey? & d \\
$\alpha_2$ & Awareness-Fit & What barrier? (cognitive/emotional/habitual) & d \\
$\alpha_3$ & Dimension-Fit & Which lever? (Info/Default/Incentive/Feedback) & pp \\
$\alpha_4$ & Scale-Fit & Which level? (Individual/Org/System) & d \\
$\alpha_5$ & Resource-Fit & Is it feasible? (Budget/Time/Implementation) & d \\
$\alpha_6$ & Context-Fit & Does it fit here? (Culture/Institution) & d \\
\bottomrule
\end{tabular}
\caption{The six $\alpha$-dimensions with their measurement scales.}
\label{tab:alpha-dimensions}
\end{table}

% -----------------------------------------------------------------------------
% SECTION: THEORY
% -----------------------------------------------------------------------------

\section{Core Theory: The LLMMC Approach}
\label{sec:iii-theory}

The LLMMC (LLM Monte Carlo) approach is based on the observation that Large Language
Models encode substantial knowledge about behavioral interventions from the scientific
literature. By systematically eliciting this knowledge and calibrating it against
empirical anchors, we can generate reliable predictions for new contexts.

\textbf{Theoretical Foundation:}
\begin{enumerate}
\item \textbf{LLM as Knowledge Aggregator}: LLMs compress millions of papers into
accessible estimates. This makes meta-analytic knowledge available for any context.
\item \textbf{Calibration as Bias Correction}: Raw LLM estimates are systematically
biased. Linear calibration corrects this bias while preserving relative ordering.
\item \textbf{Uncertainty Propagation}: Monte Carlo simulation propagates elicitation
uncertainty through the entire pipeline, ensuring honest confidence intervals.
\end{enumerate}

% -----------------------------------------------------------------------------
% SECTION: CALIBRATION
% -----------------------------------------------------------------------------

\section{Calibration Methodology}
\label{sec:iii-calibration}

\subsection{The Calibration Axioms}

\begin{tcolorbox}[colback=red!5!white, colframe=red!75!black,
    title=Axiom CAL-1: Scale-Local Calibration]
d/SMD and percentage-point (pp) scales are calibrated \textbf{separately}.
No forced integration across scales.
\end{tcolorbox}

\begin{tcolorbox}[colback=red!5!white, colframe=red!75!black,
    title=Axiom CAL-2: Empirical Anchoring]
Calibration parameters $(a, b)$ are derived from Tier-1/2 empirical anchors only.
\end{tcolorbox}

\begin{tcolorbox}[colback=red!5!white, colframe=red!75!black,
    title=Axiom CAL-3: Uncertainty Propagation]
Elicitation uncertainty (EU) is propagated through all pipeline stages via Monte Carlo.
\end{tcolorbox}

\begin{tcolorbox}[colback=red!5!white, colframe=red!75!black,
    title=Axiom CAL-4: Falsifiability]
The calibration is falsifiable: if predicted effects deviate systematically from
observed effects, the calibration must be revised.
\end{tcolorbox}

\subsection{Calibration Formulas}

\textbf{d/SMD Scale:}
\begin{equation}
\theta_{\text{true}} = 0.03 + 0.79 \times \theta_{\text{llm}} \quad (\sigma_{\text{model}} = 0.06)
\end{equation}

\textbf{Percentage-Point Scale:}
\begin{equation}
\text{pp}_{\text{true}} = -0.48 + 1.007 \times \text{pp}_{\text{llm}} \quad (\sigma_{\text{model}} = 3.54)
\end{equation}

\subsection{The 5-Point Deployment Checklist}
\label{sec:iii-checklist}

Before deploying a calibration, all five checks must pass:

\begin{table}[h]
\centering
\begin{tabular}{clcc}
\toprule
\textbf{\#} & \textbf{Check} & \textbf{Criterion} & \textbf{Status} \\
\midrule
1 & Support Coverage & $\text{pp}_{\min} \leq 10$, $\text{pp}_{\max} \geq 60$ & \checkmark \\
2 & Intercept Size & $|a| < 5\text{pp}$ & \checkmark \\
3 & Slope & $b \in [0.85, 1.05]$ & \checkmark \\
4 & LOO Coverage & $\geq 90\%$ for 95\%-CI & \checkmark \\
5 & Domain Robustness & Mean residual $< 10\text{pp}$ & \checkmark \\
\bottomrule
\end{tabular}
\caption{5-Point Deployment Checklist (all checks passed for pp calibration).}
\label{tab:checklist}
\end{table}

% -----------------------------------------------------------------------------
% SECTION: MAPPING
% -----------------------------------------------------------------------------

\section{Piecewise Linear Mapping}
\label{sec:iii-mapping}

\begin{tcolorbox}[colback=blue!5!white, colframe=blue!75!black,
    title=Axiom MAP-1: Auditable Mapping]
The mapping from effect scales to $\theta \in [0,1]$ uses piecewise linear functions
with explicit, auditable knot points. Version: \texttt{mapping\_v1}.
\end{tcolorbox}

\subsection{d/SMD Mapping Knots}

\begin{table}[h]
\centering
\begin{tabular}{cc}
\toprule
\textbf{d (calibrated)} & \textbf{$\theta$} \\
\midrule
0.0 & 0.00 \\
0.2 & 0.30 \\
0.5 & 0.60 \\
0.8 & 0.90 \\
1.2 & 0.97 \\
\bottomrule
\end{tabular}
\caption{d/SMD to $\theta$ mapping knots.}
\end{table}

\subsection{Percentage-Point Mapping Knots}

\begin{table}[h]
\centering
\begin{tabular}{cc}
\toprule
\textbf{pp (calibrated)} & \textbf{$\theta$} \\
\midrule
0 & 0.00 \\
5 & 0.30 \\
15 & 0.50 \\
30 & 0.70 \\
60 & 0.90 \\
90 & 0.97 \\
\bottomrule
\end{tabular}
\caption{Percentage-point to $\theta$ mapping knots.}
\end{table}

% -----------------------------------------------------------------------------
% SECTION: R-SCORE
% -----------------------------------------------------------------------------

\section{R-Score Calculation}
\label{sec:iii-rscore}

\begin{tcolorbox}[colback=blue!5!white, colframe=blue!75!black,
    title=Axiom RSC-1: R-Score Formula]
\begin{equation}
R(a,K) = \sum_{i=1}^{6} \alpha_i + \gamma \cdot \|a\| \cdot \|K\|
\end{equation}
where:
\begin{itemize}[nosep]
\item $\alpha_i$: Calibrated fit parameter for dimension $i$
\item $\gamma$: Complementarity coefficient (intervention $\times$ context)
\item $\|a\|$: Design norm (intervention strength)
\item $\|K\|$: Context norm (context receptivity)
\end{itemize}
\end{tcolorbox}

\textbf{Important:} R is a \emph{relevance/prioritization score}, not a causal effect.

\subsection{Monte Carlo Uncertainty Propagation}

For each parameter, draw from its uncertainty distribution:
\begin{align}
\alpha_i &\sim \mathcal{N}(\hat{\alpha}_i, \text{EU}_i) \\
\gamma &\sim \mathcal{N}(\hat{\gamma}, \text{EU}_\gamma)
\end{align}

Compute $R$ for $n = 50{,}000$ draws to obtain:
\begin{itemize}[nosep]
\item $E[R]$: Expected R-Score
\item 95\%-CI: $[R_{2.5\%}, R_{97.5\%}]$
\item $P(R > T)$: Probability of exceeding threshold
\end{itemize}

% -----------------------------------------------------------------------------
% SECTION: DECISION RULES
% -----------------------------------------------------------------------------

\section{Decision Rules}
\label{sec:iii-decision}

\begin{tcolorbox}[colback=green!5!white, colframe=green!75!black,
    title=Axiom DEC-1: Decision Rules]

Given threshold $T$ (default: 6.0):

\begin{center}
\begin{tabular}{lcl}
$P(R > T) < 5\%$ & $\Rightarrow$ & \textbf{NO-GO} \\
$5\% \leq P(R > T) < 25\%$ & $\Rightarrow$ & \textbf{PILOT} \\
$P(R > T) \geq 25\%$ & $\Rightarrow$ & \textbf{GO} \\
\end{tabular}
\end{center}

\end{tcolorbox}

% -----------------------------------------------------------------------------
% SECTION: KEY RESULTS
% -----------------------------------------------------------------------------

\section{Key Results}
\label{sec:iii-results}

\subsection{Calibration Validation}

Both calibration paths have passed the 5-point deployment checklist:

\begin{table}[h]
\centering
\begin{tabular}{lcc}
\toprule
\textbf{Metric} & \textbf{d/SMD} & \textbf{pp} \\
\midrule
Support Coverage & 0.15--0.79 & 4--87 pp \\
Intercept & 0.03 & -0.48 pp \\
Slope & 0.79 & 1.007 \\
LOO Coverage & 92\% & 100\% \\
Status & \checkmark Production & \checkmark Production \\
\bottomrule
\end{tabular}
\caption{Calibration validation results for both scales.}
\end{table}

\subsection{Predictive Accuracy}

The pipeline correctly predicts effect magnitudes across diverse contexts:
\begin{itemize}[nosep]
\item \textbf{High-$\alpha$ contexts} (e.g., defaults for low-involvement goods): 60--90 pp
\item \textbf{Medium-$\alpha$ contexts} (e.g., social norm interventions): 10--30 pp
\item \textbf{Low-$\alpha$ contexts} (e.g., information-only in action phase): 2--8 pp
\end{itemize}


% -----------------------------------------------------------------------------
% SECTION: VALIDATION - POPULATION STATISTICAL REALISM
% -----------------------------------------------------------------------------

\section{LLMMC Validation: Population Statistical Realism}
\label{sec:iii-validation-psr}

\subsection{The Typological Bias Problem}

LLMs are trained to generate plausible individual responses, not to reproduce
population-level statistical structures. This creates a fundamental validity concern
for LLMMC estimates \citep{PAP-xie2013population}.

\begin{tcolorbox}[colback=red!5!white, colframe=red!75!black,
    title=Warning: Typological Bias in LLMs]
Yu Xie (2013, 2026) identifies that LLMs exhibit \textbf{typological bias}:
\begin{itemize}[nosep]
\item Distributions \textbf{collapse} toward stereotypical values
\item Associations are \textbf{stereotyped} rather than empirically grounded
\item Life course sequences are \textbf{poorly reproduced}
\item Model size $\neq$ statistical quality
\end{itemize}

\textbf{Implication:} Individual narrative plausibility is \emph{not} a valid criterion
for LLMMC calibration. We must validate against \textbf{population-level statistical patterns}.
\end{tcolorbox}

\subsection{The SSDataBench Framework}

Following Xie et al.'s SSDataBench methodology, LLMMC outputs should be validated
against five types of statistical patterns:

\begin{tcolorbox}[colback=blue!5!white, colframe=blue!75!black,
    title=Axiom CAL-5: Population Statistical Realism (Xie Criterion)]
LLMMC estimates are valid if and only if they reproduce \textbf{population-level
statistical patterns}, not merely individual-case plausibility.

\textbf{Validation requires checking:}
\begin{enumerate}[nosep]
\item Univariate distributions (marginals)
\item Bivariate associations (correlations)
\item Multivariate outcome predictions
\item Life event sequences (for longitudinal estimates)
\item Covariate-outcome relationships
\end{enumerate}
\end{tcolorbox}

\begin{table}[h]
\centering
\caption{SSDataBench Validation Criteria for LLMMC}
\label{tab:ssdatabench}
\begin{tabular}{clp{5cm}l}
\toprule
\textbf{\#} & \textbf{Pattern Type} & \textbf{Validation Question} & \textbf{Test} \\
\midrule
1 & Univariate & Does the distribution match empirical data? & KS-test \\
2 & Bivariate & Are associations correctly signed and sized? & Correlation \\
3 & Multivariate & Can we predict outcomes from covariates? & $R^2$, RMSE \\
4 & Sequences & Are life course patterns realistic? & Sequence alignment \\
5 & Heterogeneity & Does HTE vary correctly by segment? & Subgroup $\Delta$ \\
\bottomrule
\end{tabular}
\end{table}

\subsection{Population Thinking vs. Typological Thinking}

Xie (2013) distinguishes two epistemological traditions with direct implications for LLMMC:

\begin{table}[h]
\centering
\caption{Epistemological Foundations: Typological vs. Population Thinking}
\label{tab:epistemology}
\begin{tabular}{lll}
\toprule
\textbf{Aspect} & \textbf{Typological (Physics)} & \textbf{Population (Social Science)} \\
\midrule
Model & $y_i = m + e_i$ & $y_i = m_i + e_i$ \\
Variation & Error to eliminate & Phenomenon to study \\
Mean & ``True'' parameter & Aggregate of heterogeneous units \\
Causation & Scientific solution & Statistical solution only \\
LLM Risk & Generates ``typical'' cases & Should generate distributions \\
\bottomrule
\end{tabular}
\end{table}

\textbf{Key insight from Holland (1986):}
\begin{quote}
\emph{``In social science, individual causal effects are not identifiable.
Causal inference is always group-based and assumption-specific.''}
\end{quote}

This validates EBF's context-first approach: Parameters like $\lambda$ (loss aversion)
are not fixed constants but emerge from the 5-level context cascade
(MACRO $\to$ MESO $\to$ MICRO $\to$ INDIVIDUAL $\to$ META).

\subsection{LLMMC Validation Protocol}

Based on the SSDataBench framework, we implement the following validation protocol:

\begin{tcolorbox}[colback=green!5!white, colframe=green!75!black,
    title=LLMMC Validation Protocol (Population Realism)]

\textbf{Before deploying LLMMC estimates:}

\begin{enumerate}
\item \textbf{Distribution Check:} Compare LLMMC-generated parameter distributions
against empirical distributions from Tier-1/2 studies.
\begin{itemize}[nosep]
\item Test: Kolmogorov-Smirnov $p > 0.05$
\item Red flag: Collapsed distributions (low variance)
\end{itemize}

\item \textbf{Association Check:} Verify that correlations between parameters
match empirical patterns.
\begin{itemize}[nosep]
\item Test: $|\rho_{\text{LLMMC}} - \rho_{\text{empirical}}| < 0.15$
\item Red flag: Stereotyped associations (e.g., all positive)
\end{itemize}

\item \textbf{Heterogeneity Check:} Confirm that segment-specific estimates
vary appropriately (not collapsed to mean).
\begin{itemize}[nosep]
\item Test: $\sigma_{\text{between-segment}} > 0.1 \times \bar{\theta}$
\item Red flag: All segments get same estimate
\end{itemize}

\item \textbf{Sequence Check (if applicable):} For BCJ trajectory estimates,
verify life course patterns are realistic.
\begin{itemize}[nosep]
\item Test: Sequence alignment score $> 0.7$
\item Red flag: Implausible transitions
\end{itemize}
\end{enumerate}

\textbf{Documentation:} Log all validation results in \texttt{data/llmmc\_validation\_log.yaml}

\end{tcolorbox}

\subsection{Implications for EBF Methodology}

\begin{enumerate}
\item \textbf{Calibration anchors must be population-representative:}
Tier-1/2 anchors should come from representative samples, not convenience samples.

\item \textbf{Heterogeneous treatment effects (HTE) are the norm:}
Following Brand \& Xie (2010), we expect different segments to respond differently.
The segment-multiplier framework (Chapter 19) operationalizes this.

\item \textbf{Context-specific parameters are essential:}
There is no ``true'' $\lambda$ --- only $\lambda_{i}$ conditional on context $\Psi$.
This validates the 5-level context cascade in CLAUDE.md.

\item \textbf{LLMMC is a prior, not a posterior:}
LLMMC generates informed priors that must be updated with empirical evidence
via Bayesian updating (see Section~\ref{sec:iii-calibration}).
\end{enumerate}

\textbf{References for this section:}
\begin{itemize}[nosep]
\item \citet{PAP-xie2013population}: Population Heterogeneity and Causal Inference
\item \citet{PAP-holland1986statistics}: Statistics and Causal Inference
\item \citet{PAP-brand2010benefits}: Heterogeneous Returns to Education
\item \citet{PAP-xie2012estimating}: Estimating HTE with Observational Data
\end{itemize}


% -----------------------------------------------------------------------------
% SECTION: ADAPTIVE EXPERIMENTS AND POLICY LEARNING
% -----------------------------------------------------------------------------

\section{Adaptive Experiments and Policy Learning}
\label{sec:iii-adaptive}

While Section~\ref{sec:iii-validation-psr} addresses \emph{validation} of LLMMC estimates,
this section addresses \emph{calibration improvement} through adaptive experimental designs
\citep{PAP-kasy2021adaptive, PAP-caria2024refugees}.

\subsection{The Welfare vs. Reward Distinction}

A fundamental insight from Kasy (2025, 2026) is that \textbf{welfare $\neq$ reward}:

\begin{tcolorbox}[colback=blue!5!white, colframe=blue!75!black,
    title=Axiom CAL-6: Welfare-Reward Distinction (Kasy Criterion)]
Policy learning must optimize \textbf{social welfare}, not individual rewards.
This distinction matters because:
\begin{enumerate}[nosep]
\item Multiple agents have conflicting interests
\item Distribution of outcomes matters (inequality aversion)
\item Externalities affect non-participants
\item Long-term consequences differ from short-term rewards
\end{enumerate}

\textbf{Implication for LLMMC:} Calibration should weight effects by welfare impact,
not just average treatment effects (ATE).
\end{tcolorbox}

\begin{table}[h]
\centering
\caption{Reward vs. Welfare in Policy Learning}
\label{tab:reward-welfare}
\begin{tabular}{lll}
\toprule
\textbf{Aspect} & \textbf{Reward (RL)} & \textbf{Welfare (Economics)} \\
\midrule
Objective & Maximize $\sum r_t$ & Maximize $W = \sum u_i(x)$ \\
Agents & Single agent & Multiple agents \\
Distribution & Irrelevant & Central concern \\
Externalities & Ignored & Internalized \\
Time horizon & Discount $\gamma$ & Social discount rate \\
EBF mapping & Raw effect $\Delta$ & Welfare-weighted $\alpha$ \\
\bottomrule
\end{tabular}
\end{table}

\subsection{Adaptive Experiments for Calibration}

Standard LLMMC calibration uses fixed RCT evidence. Adaptive experiments can
\emph{dynamically improve} calibration by learning which contexts and interventions
need more precise estimates \citep{PAP-kasy2021adaptive}.

\begin{tcolorbox}[colback=green!5!white, colframe=green!75!black,
    title=Exploration Sampling Algorithm (Kasy \& Sautmann 2021)]

\textbf{Problem:} Given $K$ treatments and $T$ experimental waves, assign units
to treatments to maximize policy value at the end.

\textbf{Solution:} Exploration Sampling assigns treatment $k$ with probability:
\[
p_k^{(t)} \propto f\left( P(\text{treatment } k \text{ is optimal} \mid \text{data}_{1:t-1}) \right)
\]
where $f$ is an increasing concave function (e.g., $f(x) = \sqrt{x}$).

\textbf{Properties:}
\begin{itemize}[nosep]
\item Asymptotically optimal for policy choice
\item Dominates both non-adaptive designs and standard bandits
\item Balances exploration (learning) and exploitation (welfare)
\end{itemize}

\textbf{For LLMMC:} Use exploration sampling to allocate calibration effort
across intervention types, prioritizing high-uncertainty, high-impact domains.
\end{tcolorbox}

\subsection{Mechanism Design as Causal Learning}

Kasy (2026) argues that \textbf{mechanism design is causal inference}:

\begin{itemize}
\item \textbf{VCG mechanisms} (Vickrey-Clarke-Groves) elicit private information
through incentive-compatible payments \citep{PAP-vickrey1961counterspeculation,
PAP-clarke1971multipart, PAP-groves1973incentives}

\item \textbf{Pigouvian taxation} reveals social costs through optimal corrective
taxes \citep{PAP-pigou1920welfare}

\item \textbf{BDM mechanism} elicits willingness-to-pay through incentive-compatible
lotteries \citep{PAP-becker1964bdm}
\end{itemize}

\begin{tcolorbox}[colback=yellow!5!white, colframe=yellow!75!black,
    title=Key Insight: Causality $\Leftrightarrow$ Decision]

Kasy's central argument:
\begin{quote}
\emph{``Reinforcement learning is inherently causal. Every intervention is an experiment.
Mechanism design learns causal parameters through incentive-compatible elicitation.''}
\end{quote}

\textbf{Implication for EBF:} The 9D intervention vector $\vec{I} \in [0,1]^9$
(Appendix IE) can be learned through adaptive mechanisms, not just estimated from
historical data.
\end{tcolorbox}

\subsection{Double/Debiased ML for LLMMC}

When LLMMC estimates depend on high-dimensional covariates, the Double/Debiased ML
framework \citep{PAP-chernozhukov2018double} provides valid inference:

\begin{enumerate}
\item \textbf{Nuisance estimation:} Use ML to estimate propensity scores and
outcome models (first stage)

\item \textbf{Cross-fitting:} Split sample into folds; estimate nuisance on
one fold, target parameter on others

\item \textbf{Neyman orthogonality:} Construct moment conditions that are
insensitive to nuisance parameter estimation error

\item \textbf{Result:} Valid confidence intervals for treatment effects despite
high-dimensional controls
\end{enumerate}

\textbf{For LLMMC:} When calibrating context-specific parameters $\theta(\Psi)$,
use DML to control for high-dimensional context features without overfitting.

\subsection{Causal Forests for Heterogeneous Calibration}

Causal Forests \citep{PAP-athey2016recursive, PAP-wager2018causalforest} enable
segment-specific LLMMC calibration with valid inference:

\begin{itemize}
\item \textbf{Honest estimation:} Separate trees for structure (which covariates
matter) and effect estimation (magnitudes)

\item \textbf{Asymptotic normality:} Valid confidence intervals for conditional
average treatment effects (CATE)

\item \textbf{Heterogeneity detection:} Automatically identifies segments with
different treatment responses
\end{itemize}

This connects to the EBF segment-multiplier framework (Chapter 19): Causal Forests
can estimate segment-specific $\sigma_s$ values from RCT data, which then calibrate
LLMMC heterogeneity estimates.

\subsection{Implications for LLMMC Calibration Pipeline}

\begin{enumerate}
\item \textbf{Welfare-weighted calibration:}
Weight calibration anchors by welfare impact, not just sample size.
High-stakes domains (health, environment) deserve tighter calibration.

\item \textbf{Adaptive calibration allocation:}
Use exploration sampling to allocate calibration effort across intervention
types and contexts. Prioritize high-uncertainty, high-welfare-impact cells.

\item \textbf{DML for context-rich estimation:}
When context $\Psi$ is high-dimensional, use DML to estimate $\alpha(\Psi)$
without overfitting to context features.

\item \textbf{Causal Forests for segmentation:}
Use Causal Forests to identify segments with heterogeneous treatment effects,
then calibrate LLMMC segment multipliers to match.

\item \textbf{Mechanism-based validation:}
Consider incentive-compatible elicitation (BDM, prediction markets) as
alternative validation for LLMMC-generated WTP estimates.
\end{enumerate}

\textbf{References for this section:}
\begin{itemize}[nosep]
\item \citet{PAP-kasy2021adaptive}: Adaptive Treatment Assignment (Econometrica 2021)
\item \citet{PAP-kasy2025adaptive}: Adaptive Welfare Maximization (Econometrica 2025)
\item \citet{PAP-chernozhukov2018double}: Double/Debiased ML (Econometrics Journal 2018)
\item \citet{PAP-athey2016recursive}: Causal Trees (PNAS 2016)
\item \citet{PAP-wager2018causalforest}: Causal Forests (JASA 2018)
\item \citet{PAP-caria2024refugees}: Adaptive Field Experiment (JEEA 2024)
\item \citet{PAP-becker1964bdm}: BDM Mechanism (Behavioral Science 1964)
\end{itemize}


% -----------------------------------------------------------------------------
% SECTION: CAUSES OF EFFECTS AND RETROSPECTIVE ATTRIBUTION
% -----------------------------------------------------------------------------

\section{Causes of Effects: Retrospective Individual Attribution}
\label{sec:iii-coe}

While Sections~\ref{sec:iii-validation-psr} and \ref{sec:iii-adaptive} address
\emph{effects of causes} (forward-looking, population-level), this section addresses
the complementary problem of \emph{causes of effects}---retrospective, individual-level
causal attribution \citep{PAP-pearl2015causes, PAP-dawid2022effects}.

\subsection{The Fundamental Distinction}

\begin{tcolorbox}[colback=red!5!white, colframe=red!75!black,
    title=Effects of Causes vs. Causes of Effects]

\textbf{Effects of Causes (EoC):}
Given a cause $X$, what is its effect on $Y$?
\begin{itemize}[nosep]
\item Forward-looking, predictive
\item Population-level (ATE, CATE)
\item Standard causal inference
\end{itemize}

\textbf{Causes of Effects (CoE):}
Given an observed effect $Y=y$, what caused it?
\begin{itemize}[nosep]
\item Backward-looking, retrospective
\item Individual-level attribution
\item Legal and moral reasoning
\end{itemize}

\textbf{Key insight:} EoC conditions on $X$; CoE conditions on both $X$ \emph{and} $Y$.
\end{tcolorbox}

\subsection{Probability of Causation}

The \textbf{Probability of Causation (PoC)} framework \citep{PAP-tian2000probabilities}
provides three key estimands for individual attribution:

\begin{table}[h]
\centering
\caption{Probability of Causation Estimands}
\label{tab:poc}
\begin{tabular}{lll}
\toprule
\textbf{Estimand} & \textbf{Definition} & \textbf{Legal Interpretation} \\
\midrule
PN (Necessity) & $P(Y_0 = 0 \mid X = 1, Y = 1)$ & ``But-for'' causation \\
PS (Sufficiency) & $P(Y_1 = 1 \mid X = 0, Y = 0)$ & Would treatment have helped? \\
PNS (Both) & $P(Y_1 = 1, Y_0 = 0)$ & Necessary \emph{and} sufficient \\
\bottomrule
\end{tabular}
\end{table}

\textbf{PN} asks: Given that treatment was applied ($X=1$) and the outcome occurred ($Y=1$),
what is the probability that the outcome would \emph{not} have occurred without treatment?
This corresponds to the legal ``but-for'' test.

\subsection{Bounds and Identification}

PoC estimands are not point-identified from observational data alone, but
\textbf{informative bounds} can be derived \citep{PAP-tian2000probabilities}:

\begin{tcolorbox}[colback=green!5!white, colframe=green!75!black,
    title=Tian-Pearl Bounds on PN]

From observational data $(P(Y|X))$ and experimental data $(P(Y|do(X)))$:
\[
\max\left(0, \frac{P(Y=1|X=1) - P(Y=1|do(X=0))}{P(Y=1|X=1)}\right) \leq PN
\]
\[
PN \leq \min\left(1, \frac{P(Y=0|do(X=0))}{P(Y=1|X=1)}\right)
\]

\textbf{Key insight:} Fusion of observational and experimental data tightens bounds.
\end{tcolorbox}

\subsection{Posterior Causal Estimands}

Recent work extends PoC to continuous outcomes \citep{PAP-luo2025attribution}:

\begin{tcolorbox}[colback=blue!5!white, colframe=blue!75!black,
    title=Axiom CAL-7: Posterior Causal Attribution (Luo-Geng Criterion)]
For individual causal attribution with continuous outcomes, condition on
\textbf{both} the treatment received and the outcome observed:
\[
\theta^{PCE} = E[Y(0) \mid X = 1, Y = y]
\]
This ``posterior causal estimand'' asks: What would this individual's outcome
have been under no treatment, given that they received treatment and experienced $Y=y$?
\end{tcolorbox}

\subsection{Implications for LLMMC and EBF}

\begin{enumerate}
\item \textbf{Individual Attribution:}
LLMMC typically estimates population-level effects (EoC). For individual attribution
(e.g., ``Did this intervention cause this person's behavior change?''), PoC bounds
provide a principled framework.

\item \textbf{Legal and Regulatory Applications:}
When EBF analyses inform legal or regulatory decisions, PN bounds can translate
statistical evidence into legally relevant probabilities \citep{PAP-greenland2000epidemiology}.

\item \textbf{Segment-to-Individual Inference:}
Dawid and Musio's work on ``group to individual'' inference \citep{PAP-dawid2022individual}
provides bounds for attributing segment-level effects to individuals---key for
EBF's segment framework (Chapter 19).

\item \textbf{Complementarity of Statistics and Decisions:}
Statistics quantifies causal contribution under uncertainty; normative frameworks
(legal standards, ethical principles) determine acceptable thresholds for action.
This separation is essential for responsible use of LLMMC in high-stakes contexts.
\end{enumerate}

\textbf{References for this section:}
\begin{itemize}[nosep]
\item \citet{PAP-pearl2015causes}: Causes of Effects and Effects of Causes (SMR 2015)
\item \citet{PAP-dawid2022effects}: Effects of Causes and Causes of Effects (Ann Rev Stats 2022)
\item \citet{PAP-tian2000probabilities}: Probabilities of Causation: Bounds and Identification
\item \citet{PAP-luo2025attribution}: Causal Attribution for Continuous Outcomes (ICML 2025)
\item \citet{PAP-greenland2000epidemiology}: Epidemiology, Justice, and PoC (Jurimetrics 2000)
\item \citet{PAP-hart1985causation}: Causation in the Law (OUP 1985)
\end{itemize}


% -----------------------------------------------------------------------------
% SECTION: SEQUENTIAL CAUSAL INFERENCE
% -----------------------------------------------------------------------------

\section{Sequential Causal Inference and Trajectory Decomposition}
\label{sec:iii-seq}

While previous sections address single-stage interventions, many EBF applications
involve \emph{sequences} of treatments over time: customer journeys, adaptive
experiments, and dynamic treatment regimes. This section integrates sequential
causal inference methods into the LLMMC framework, drawing on recent work by
Syrgkanis and Zhan \citep{PAP-syrgkanis2025postrl}.

\subsection{From Single-Stage to Sequential Effects}

\begin{tcolorbox}[colback=red!5!white, colframe=red!75!black,
    title=The Sequential Challenge]

\textbf{Single-stage:} One treatment $A$, one outcome $Y$.
\[
\tau = E[Y(1) - Y(0)]
\]

\textbf{Sequential:} Treatments $A_1, A_2, \ldots, A_T$ over time, outcome $Y$.
\begin{itemize}[nosep]
\item Each $A_t$ depends on history $H_t$
\item Effects compound and interact
\item Standard methods (ATE, CATE) don't decompose well
\end{itemize}

\textbf{Key question:} How much of the total effect is attributable to each stage?
\end{tcolorbox}

\subsection{Blip Functions: Stage-Wise Causal Effects}

The \textbf{blip function} \citep{PAP-robins2000msm, PAP-syrgkanis2025postrl}
isolates the causal effect of treatment at stage $t$, conditional on history:

\begin{tcolorbox}[colback=blue!5!white, colframe=blue!75!black,
    title=Axiom CAL-8: Blip Decomposition (Syrgkanis-Zhan)]
The stage-$t$ blip function decomposes the total trajectory effect:
\[
\gamma_t(H_t, A_t) = E[Y^* - Y^{A_t \leftarrow 0} \mid H_t, A_t]
\]
where $Y^*$ is the observed outcome and $Y^{A_t \leftarrow 0}$ is the counterfactual
outcome had treatment at stage $t$ been set to zero (following optimal thereafter).

\textbf{Total effect decomposition:}
\[
E[Y^{d^*}] - E[Y^{d_0}] = \sum_{t=1}^{T} E[\gamma_t(H_t, d^*_t(H_t))]
\]
where $d^*$ is the optimal regime and $d_0$ is the baseline (no treatment) regime.
\end{tcolorbox}

\textbf{Interpretation:} The blip function answers: ``What is the causal effect of
the treatment decision at stage $t$, for someone with history $H_t$?'' This is
precisely what we need for customer journey analysis in EBF---understanding which
touchpoints matter most, conditional on journey history.

\subsection{Dynamic Treatment Regimes}

A \textbf{dynamic treatment regime} (DTR) \citep{PAP-murphy2003optimal} is a sequence
of decision rules $d = (d_1, d_2, \ldots, d_T)$ where each $d_t(H_t)$ maps history
to treatment:

\begin{itemize}
\item \textbf{Value function:} $V^d = E[Y \mid \text{follow regime } d]$
\item \textbf{Optimal regime:} $d^* = \arg\max_d V^d$
\item \textbf{Q-learning:} Backward induction starting from $Q_T$, then $Q_{T-1}$, etc.
\end{itemize}

\textbf{For EBF:} Customer journey optimization is DTR. The optimal ``nudge sequence''
for a customer depends on their journey history---their current stage, past responses,
and accumulated engagement.

\subsection{Marginal Structural Models for Time-Varying Confounding}

When treatments are observational (not randomized), \textbf{Marginal Structural Models}
(MSM) \citep{PAP-robins2000msm} handle time-varying confounders:

\begin{enumerate}
\item \textbf{Problem:} Confounder $L_t$ affects both treatment $A_t$ and outcome $Y$,
\emph{and} $L_t$ is affected by prior treatment $A_{t-1}$.

\item \textbf{Solution:} Inverse Probability of Treatment Weighting (IPTW):
\[
W_i = \prod_{t=1}^{T} \frac{1}{P(A_{it} \mid H_{it})}
\]

\item \textbf{Stabilized weights:} Reduce variance by including numerator:
\[
SW_i = \prod_{t=1}^{T} \frac{P(A_{it} \mid A_{i,t-1})}{P(A_{it} \mid H_{it})}
\]
\end{enumerate}

\textbf{For LLMMC:} When calibrating against observational journey data, MSM weights
correct for selection bias from confounded treatment sequences.

\subsection{Trajectory Decomposition: Extending Oaxaca-Blinder}

The classic \textbf{Oaxaca-Blinder decomposition} \citep{PAP-oaxaca1973decomposition}
decomposes group differences:
\[
\bar{Y}_A - \bar{Y}_B = \underbrace{(\bar{X}_A - \bar{X}_B)'\hat{\beta}}_{\text{Explained}} + \underbrace{\bar{X}_B'(\hat{\beta}_A - \hat{\beta}_B)}_{\text{Unexplained}}
\]

\textbf{For sequences:} Zhan's work extends this to trajectories:

\begin{tcolorbox}[colback=green!5!white, colframe=green!75!black,
    title=Trajectory Gap Decomposition]
The gap in outcomes between groups $A$ and $B$ with different treatment trajectories:
\[
E[Y_A] - E[Y_B] = \underbrace{\sum_t E[\gamma_t](\pi_t^A - \pi_t^B)}_{\text{Distribution Effect}} + \underbrace{\sum_t (\gamma_t^A - \gamma_t^B)E[\pi_t]}_{\text{Blip Effect}}
\]
where $\pi_t$ is the probability of receiving treatment at stage $t$.

\textbf{Interpretation:}
\begin{itemize}[nosep]
\item \textbf{Distribution Effect:} Group differences due to different treatment \emph{rates}
at each stage (e.g., group A gets more touchpoints)
\item \textbf{Blip Effect:} Group differences due to different treatment \emph{effects}
at each stage (e.g., group A responds more to each touchpoint)
\end{itemize}
\end{tcolorbox}

\subsection{Implications for LLMMC and BCJ}

\begin{enumerate}
\item \textbf{Stage-wise LLMMC calibration:}
Instead of calibrating total journey effects, calibrate blip functions at each stage.
This enables more granular intervention design.

\item \textbf{Customer Journey Modeling (BCJ):}
The EBF Behavioral Customer Journey framework (Chapter 13) can be formalized using
DTR. LLMMC can estimate stage-specific treatment effects $\gamma_t$ conditional on
journey history $H_t$.

\item \textbf{Journey gap analysis:}
Use trajectory decomposition to explain why different customer segments have different
outcomes. Is it because they receive different touchpoints (distribution) or respond
differently to the same touchpoints (blip)?

\item \textbf{Optimal journey design:}
Q-learning with LLMMC-estimated blips can identify the optimal nudge sequence for
each customer segment.

\item \textbf{RL-generated data calibration:}
When LLMMC estimates are calibrated against data from adaptive experiments or
recommendation systems, use post-RL inference methods \citep{PAP-syrgkanis2025postrl}
to account for non-stationarity.
\end{enumerate}

\textbf{References for this section:}
\begin{itemize}[nosep]
\item \citet{PAP-syrgkanis2025postrl}: Post RL Inference (Operations Research 2025)
\item \citet{PAP-zhang2024personalized}: Personalized Policy Learning (Management Science 2024)
\item \citet{PAP-robins2000msm}: Marginal Structural Models (Epidemiology 2000)
\item \citet{PAP-murphy2003optimal}: Optimal Dynamic Treatment Regimes (JRSS-B 2003)
\item \citet{PAP-oaxaca1973decomposition}: Male-Female Wage Differentials (IER 1973)
\item \citet{PAP-athey2021policy}: Policy Learning with Observational Data (ECMA 2021)
\end{itemize}


% -----------------------------------------------------------------------------
% SECTION: CIFW05 WORKSHOP ROADMAP
% -----------------------------------------------------------------------------

\section{CIFW05 Workshop: Conceptual Map and Integration Roadmap}
\label{sec:iii-cifw05-roadmap}

The CIFW05 Workshop ``Causal Machine Learning for the Social Sciences'' (Newton Institute,
January 2026) can be read as a movement along three axes:

\begin{enumerate}
\item \textbf{What is causality?} (Conceptual / normative)
\item \textbf{How is heterogeneity modeled causally?} (Individuals, subgroups, populations)
\item \textbf{How does causality unfold over time \& sequences?} (Trajectories)
\end{enumerate}

\subsection{Central Thesis of the Workshop}

\begin{tcolorbox}[colback=yellow!5!white, colframe=yellow!75!black,
    title=CIFW05 Central Thesis]
\textbf{Causality in the social sciences is shifting}\\[0.5em]
\emph{from} average, static effects\\
\emph{to} individual, dynamic, explanatory trajectories.

\begin{itemize}[nosep]
\item \textbf{Deakin}: Why this matters normatively
\item \textbf{Kasy}: Why this is not value-neutral
\item \textbf{Xie / Kennedy}: How heterogeneity is formalized
\item \textbf{Athey / Zhan}: How processes are causally reconstructed
\item \textbf{Newey / Chernozhukov}: How this is identified \& estimated
\end{itemize}
\end{tcolorbox}

\subsection{Integration Status}

\textbf{Integrated into this Appendix:}

\begin{table}[h]
\centering
\begin{tabular}{llll}
\toprule
\textbf{Presenter} & \textbf{Topic} & \textbf{Section} & \textbf{Axiom} \\
\midrule
Yu Xie & Population Heterogeneity & \S\ref{sec:iii-validation-psr} & CAL-5 \\
Maximilian Kasy & ML, Causal Inference, Economics & \S\ref{sec:iii-adaptive} & CAL-6 \\
Simon Deakin (Essay) & Causation in Law \& Statistics & \S\ref{sec:iii-coe} & CAL-7 \\
Susan Athey / Ruohan Zhan & Sequence Data \& Decompositions & \S\ref{sec:iii-seq} & CAL-8 \\
\bottomrule
\end{tabular}
\caption{CIFW05 content integrated into Appendix CAL.}
\end{table}

\subsection{Topics Pending Integration}

The following CIFW05 presentations have \textbf{not yet been integrated} and represent
a research agenda for extending LLMMC:

\begin{table}[h]
\centering
\small
\begin{tabular}{p{3cm}p{4cm}p{6cm}}
\toprule
\textbf{Presenter} & \textbf{Topic} & \textbf{Relevance for LLMMC} \\
\midrule
Edward Kennedy & Nonparametrics \& Heterogeneity &
Distributional effects, density estimation instead of means.
How to calibrate LLMMC for full effect distributions? \\[0.5em]

David DeGarmo & Differential Adherence \& Response &
Compliance $\neq$ effect. Critical for intervention
implementation: predicted vs. realized effects. \\[0.5em]

Whitney Newey & Nonseparable Panel Data &
Structural identification with panel data. Foundation for
long-term LLMMC calibration with longitudinal anchors. \\[0.5em]

Pasco Fearon & Birth Cohorts \& Developmental Trajectories &
Life-course perspective. Links to BCJ (Behavioral Customer Journey)
and intergenerational effects. \\[0.5em]

Alexandra Havdahl & Longitudinal Data \& Trajectories &
Empirical validation of process logic. Norwegian registry data
as potential calibration source. \\[0.5em]

Victor Chernozhukov & AI in Demand Analysis &
High-dimensional demand estimation. Proof that causal ML
is practically relevant for economic policy. \\
\bottomrule
\end{tabular}
\caption{CIFW05 topics pending integration into EBF/LLMMC framework.}
\label{tab:cifw05-pending}
\end{table}

\subsection{Workshop Structure: Three Axes}

\textbf{Axis I: What is Causality?} (Conceptual / Normative)
\begin{itemize}[nosep]
\item \textbf{Deakin}: Law = postdictive, individual, normative vs.
Statistics = predictive, generic, probabilistic
\item \textbf{Kasy}: ML $\neq$ neutral estimation; trade-offs between
efficiency vs. fairness, prediction vs. decision
\end{itemize}

\textbf{Axis II: Heterogeneity} (From Average Effects to Individuals)
\begin{itemize}[nosep]
\item \textbf{Xie}: Heterogeneity as structure, not noise; macro-level individualization
\item \textbf{Kennedy}: Densities instead of means; nonparametric identification
\item \textbf{DeGarmo}: Differential adherence as empirical stress-test
\end{itemize}

\textbf{Axis III: Time \& Sequences} (Causality as Process)
\begin{itemize}[nosep]
\item \textbf{Athey/Zhan}: Decomposition over time; generative models as counterfactual engines
\item \textbf{Newey}: Nonseparable effects in panel data; structural identification
\item \textbf{Fearon/Havdahl}: Life-course empirics; developmental trajectories
\item \textbf{Chernozhukov}: Practical landing zone for AI in economics
\end{itemize}

\subsection{Strategic Implications}

The workshop implies a \textbf{research agenda} for LLMMC:

\begin{quote}
\emph{``From causal effects to causal lives''}
\end{quote}

\begin{enumerate}
\item \textbf{Distributional calibration}: Move from mean effects to effect distributions
(Kennedy). LLMMC should estimate $P(\tau_i \in [a,b])$, not just $E[\tau]$.

\item \textbf{Compliance modeling}: Separate predicted vs. realized effects (DeGarmo).
LLMMC estimates $\neq$ implementation outcomes.

\item \textbf{Longitudinal anchors}: Use panel/registry data for calibration (Newey, Havdahl).
Norwegian registries, German SOEP, Swiss SHP as potential sources.

\item \textbf{Life-course integration}: Connect LLMMC to BCJ framework (Fearon).
Customer journeys as causal trajectories.

\item \textbf{Demand estimation}: Validate LLMMC against structural demand models (Chernozhukov).
High-dimensional economic applications.
\end{enumerate}


% =============================================================================
% PART 3: APPLICATION
% =============================================================================

\section{Worked Example: Green Energy Default (A10)}
\label{sec:iii-example}

\subsection{The Case}

\textbf{Intervention:} Green Energy Default (Opt-out instead of Opt-in)

\textbf{Context:} German households, energy provider

\textbf{Source:} Ebeling \& Lotz (2015), Energy Policy

\textbf{Effect:} 69pp difference (Default vs. Opt-in), $n = 41{,}952$ households

\subsection{Stage 1: The $\alpha$-Analysis (Filter)}

\begin{table}[h]
\centering
\begin{tabular}{clcp{6cm}}
\toprule
\textbf{Dim} & \textbf{Name} & \textbf{$\alpha_i$} & \textbf{Rationale} \\
\midrule
$\alpha_1$ & Phase-Fit & 0.95 & Trigger phase: clear decision point (contract) \\
$\alpha_2$ & Awareness-Fit & 0.85 & Known, but action barrier (status-quo bias) \\
$\alpha_3$ & Dimension-Fit & 0.98 & Default = ideal lever for passive decisions \\
$\alpha_4$ & Scale-Fit & 0.90 & Utility scale, automatable \\
$\alpha_5$ & Resource-Fit & 0.92 & No additional costs (IT change only) \\
$\alpha_6$ & Context-Fit & 0.88 & DACH: high acceptance for ``green'' defaults \\
\midrule
& \textbf{Total} & \textbf{5.48} & \emph{Exceptionally high fit} \\
\bottomrule
\end{tabular}
\caption{$\alpha$-Analysis for Green Energy Default.}
\end{table}

\textbf{Context Diagnosis ($\Psi$):}
\begin{itemize}[nosep]
\item $\Psi_1$ Low Attention: Electricity = low-involvement good
\item $\Psi_2$ Decision Cost: Tariff switch = annoying (forms, research)
\item $\Psi_3$ Trust: Utilities seen as serious, default = implicit recommendation
\end{itemize}

\textbf{Mechanism Match:} Default exploits inertia $\rightarrow$ perfect match with context.

\subsection{Stage 2: The $U$-Analysis (Utility Model)}

\textbf{The neoclassical critique:}
\begin{quote}
\emph{``If people jump from 1\% to 69\%, you manipulated their preferences!''}
\end{quote}

\textbf{Our answer:} No. The utility function $U$ was constant:
\begin{itemize}[nosep]
\item Most people have slight preference for green: $v(\text{green}) > v(\text{grey})$
\item But they also have transaction costs: $c(\text{switch}) > 0$
\end{itemize}

\textbf{The model:}
\begin{equation}
U(x) = v(x) - c(x)
\end{equation}

\begin{itemize}[nosep]
\item \textbf{Opt-in:} $U(\text{green}) - c(\text{switch}) < U(\text{grey})$ $\Rightarrow$ 1\% choose green
\item \textbf{Opt-out:} $U(\text{green}) > U(\text{grey}) - c(\text{switch})$ $\Rightarrow$ 69\% stay green
\end{itemize}

\textbf{Conclusion:} We didn't change $U$. We removed the barrier blocking $U$.

\subsection{R-Score Result}

\begin{verbatim}
E[R] = 5.7
95%-CI: [4.8, 6.6]
P(R > 5.0) = 78%

Decision: GO
\end{verbatim}

% =============================================================================
% PART 4: BACK MATTER
% =============================================================================

\section{Summary}
\label{sec:iii-summary}

\begin{tcolorbox}[colback=gray!5!white, colframe=gray!75!black,
    title=Summary: LLMMC Pipeline]

\textbf{The Pipeline:}
\begin{enumerate}[nosep]
\item LLM elicits effect estimates with structured prompts
\item Scale-specific calibration (d/SMD or pp)
\item Piecewise linear mapping to $\theta \in [0,1]$
\item R-Score computation with Monte Carlo uncertainty
\item GO/PILOT/NO-GO decision
\end{enumerate}

\textbf{Key Insight:}
\begin{quote}
\emph{``$\alpha$ is a meta-quantity over models, not a variable in models.
We make models applicable without rewriting them normatively.''}
\end{quote}

\textbf{Implementation:}
\begin{itemize}[nosep]
\item Pipeline: \texttt{scripts/llmmc\_to\_rscore.py}
\item Calibration: \texttt{data/calibration/calibration\_d\_v1.json}, \texttt{calibration\_pp\_v1.json}
\item Mapping: \texttt{scripts/theta\_mapping.py}
\item R-Score: \texttt{scripts/r\_score.py}
\end{itemize}

\end{tcolorbox}

% -----------------------------------------------------------------------------
% GLOSSARY
% -----------------------------------------------------------------------------

\section{Glossary of Symbols}
\label{sec:iii-glossary}

\begin{table}[h]
\centering
\begin{tabular}{clp{8cm}}
\toprule
\textbf{Symbol} & \textbf{Name} & \textbf{Definition} \\
\midrule
$\alpha_i$ & Alpha & Calibrated fit parameter for dimension $i$ \\
$\gamma$ & Gamma & Complementarity coefficient \\
$R$ & R-Score & Relevance/prioritization score \\
$\theta$ & Theta & Unified scale $\in [0,1]$ \\
EU & Elicitation Uncertainty & Uncertainty from LLM elicitation \\
$\|a\|$ & Design Norm & Intervention strength \\
$\|K\|$ & Context Norm & Context receptivity \\
$T$ & Threshold & Decision threshold (default: 6.0) \\
\bottomrule
\end{tabular}
\caption{Glossary of symbols. See also Appendix GLS (Central Glossary).}
\end{table}

% -----------------------------------------------------------------------------
% CRITICAL FOUNDATIONS
% -----------------------------------------------------------------------------

\section{Critical Foundations}
\label{sec:iii-foundations}

\begin{enumerate}
\item \textbf{LLM Knowledge Validity}: The approach assumes LLMs encode reliable
knowledge from the scientific literature. This is supported by extensive validation
against Tier-1/2 empirical anchors.

\item \textbf{Linear Calibration Sufficiency}: Linear calibration is sufficient
for bias correction. Non-linear patterns would require more complex calibration
functions and more anchors.

\item \textbf{Uncertainty Quantification}: Monte Carlo propagation provides honest
confidence intervals. However, this assumes uncertainties are approximately Gaussian.

\item \textbf{Context Transferability}: The calibration assumes effect magnitudes
transfer across similar contexts. Domain-specific calibrations may be needed for
highly specialized domains.

\item \textbf{Temporal Stability}: Calibration parameters are assumed stable over
time. Periodic re-validation is recommended (annual review cycle).

\item \textbf{Population Statistical Realism (Xie 2013)}: LLMs exhibit typological
bias --- they generate plausible individual cases but may not reproduce population-level
statistical patterns. LLMMC validation must check distributions, associations, and
heterogeneity against empirical benchmarks (see Section~\ref{sec:iii-validation-psr}).

\item \textbf{Welfare-Reward Distinction (Kasy 2025)}: Policy learning must optimize
social welfare, not individual rewards. LLMMC calibration should weight effects by
welfare impact, considering distributional concerns and externalities
(see Section~\ref{sec:iii-adaptive}).

\item \textbf{Effects-of-Causes vs. Causes-of-Effects (Pearl/Dawid 2015)}: LLMMC
primarily addresses effects-of-causes (forward-looking, population). For causes-of-effects
(retrospective, individual attribution), PoC bounds provide the appropriate framework.
Statistics quantifies; normative standards decide (see Section~\ref{sec:iii-coe}).

\item \textbf{Sequential Effects via Blip Decomposition (Syrgkanis/Zhan 2025)}: For
multi-stage interventions (customer journeys, adaptive experiments), blip functions
decompose total effects into stage-wise contributions. LLMMC should calibrate blips
$\gamma_t(H_t, A_t)$ conditional on history, not just aggregate trajectory effects
(see Section~\ref{sec:iii-seq}).
\end{enumerate}

% -----------------------------------------------------------------------------
% OPEN ISSUES
% -----------------------------------------------------------------------------

\section{Open Issues and Research Agenda}
\label{sec:iii-open-issues}

\begin{enumerate}
\item \textbf{Non-linear Calibration}: When are non-linear calibration functions
needed? Can we detect this from residual patterns?

\item \textbf{Domain-Specific Calibrations}: Should health, finance, and energy
domains have separate calibration parameters?

\item \textbf{LLM Model Dependence}: How stable are estimates across different
LLM versions (GPT-4, Claude, etc.)?

\item \textbf{Real-Time Updating}: Can calibrations be updated in real-time as
new RCT evidence becomes available?

\item \textbf{Heterogeneity Modeling}: How can we model heterogeneity in $\alpha$
across sub-populations?

\item \textbf{SSDataBench Implementation}: Implement the full SSDataBench validation
protocol (Section~\ref{sec:iii-validation-psr}) with automated testing against
population-level benchmarks. Priority: HIGH.

\item \textbf{Typological Bias Detection}: Develop diagnostics to detect when
LLMMC outputs exhibit typological bias (collapsed distributions, stereotyped
associations). Can we flag problematic estimates automatically?

\item \textbf{Training on Historical Data}: Xie (2026) shows that training LLMs
on historical survey data improves population realism. Can we fine-tune LLMMC
prompts with historical CPS/GSS/NLSY patterns?

\item \textbf{Adaptive Calibration Allocation}: Implement exploration sampling
(Section~\ref{sec:iii-adaptive}) to dynamically allocate calibration effort
across intervention types. Which domains need tighter calibration? Priority: MEDIUM.

\item \textbf{Welfare-Weighted Calibration}: Extend the $\alpha$-framework to
incorporate welfare weights (Kasy 2025). How should calibration prioritize
high-stakes domains (health, climate) over low-stakes domains?

\item \textbf{DML Integration}: Implement Double/Debiased ML for context-rich
LLMMC estimation. Can we estimate $\alpha(\Psi)$ with high-dimensional context
vectors without overfitting?

\item \textbf{Causal Forest Segmentation}: Use Causal Forests to discover
heterogeneous segments from RCT data, then calibrate LLMMC segment multipliers
($\sigma_s$) to match. Connect to Chapter 19 segment framework.

\item \textbf{Mechanism-Based Validation}: Explore incentive-compatible
elicitation (BDM, prediction markets) as alternative validation for LLMMC
WTP estimates. When do mechanism-elicited values diverge from survey-based values?

\item \textbf{Causes-of-Effects Bounds}: Implement PN/PS/PNS bound calculations
(Section~\ref{sec:iii-coe}) for individual attribution. Can LLMMC provide priors
for posterior causal estimands? Key for legal/medical applications.

\item \textbf{Segment-to-Individual Inference}: Develop methodology to bound
individual-level effects from segment-level LLMMC estimates. Connect to Dawid-Musio
(2022) group-to-individual framework. Critical for personalized interventions.

\item \textbf{Blip Function Calibration}: Extend LLMMC to estimate stage-wise blip
functions $\gamma_t(H_t, A_t)$ for sequential interventions. How should we elicit
conditional effects given journey history? Priority: HIGH for BCJ integration.

\item \textbf{Dynamic Treatment Regime Optimization}: Implement Q-learning with
LLMMC-estimated blips to identify optimal intervention sequences. Can LLMMC
provide sufficient precision for dynamic programming over journeys?

\item \textbf{Trajectory Decomposition Diagnostics}: Develop tools to decompose
customer segment gaps into distribution effects (who gets treated) vs. blip
effects (how they respond). Connect to Oaxaca-Blinder framework.

\item \textbf{MSM Weighting for Observational Journeys}: When calibrating against
observational customer data, implement IPTW/stabilized weights to correct for
time-varying confounding. What history variables should condition the weights?

\item \textbf{Post-RL Inference for Adaptive Systems}: When calibrating against
data from recommendation engines or adaptive experiments, use post-RL inference
methods to account for non-stationary data generating processes. Priority: MEDIUM.
\end{enumerate}

% -----------------------------------------------------------------------------
% REFERENCES
% -----------------------------------------------------------------------------

\section{References}
\label{sec:iii-references}

\textbf{Primary Sources (Intervention Effects):}
\begin{itemize}[nosep]
\item Ebeling, F., \& Lotz, S. (2015). Domestic uptake of green energy promoted by opt-out tariffs. \emph{Energy Policy}, 80, 14--22.
\item Johnson, E. J., \& Goldstein, D. (2003). Do defaults save lives? \emph{Science}, 302(5649), 1338--1339.
\item Thaler, R. H., \& Benartzi, S. (2004). Save more tomorrow. \emph{Journal of Political Economy}, 112(S1), S164--S187.
\item Hallsworth, M., et al. (2017). The behavioralist as tax collector. \emph{Journal of Public Economics}, 148, 14--31.
\end{itemize}

\textbf{Population Heterogeneity and LLMMC Validation (Section~\ref{sec:iii-validation-psr}):}
\begin{itemize}[nosep]
\item Holland, P. W. (1986). Statistics and Causal Inference. \emph{JASA}, 81(396), 945--960. [PAP-holland1986statistics]
\item Xie, Y. (2013). Population Heterogeneity and Causal Inference. \emph{PNAS}, 110(16), 6262--6268. [PAP-xie2013population]
\item Brand, J. E., \& Xie, Y. (2010). Who Benefits Most from College? \emph{ASR}, 75(2), 273--302. [PAP-brand2010benefits]
\item Xie, Y., Brand, J. E., \& Jann, B. (2012). Estimating Heterogeneous Treatment Effects. \emph{Sociological Methodology}, 42(1), 314--347. [PAP-xie2012estimating]
\item Zhou, X., \& Xie, Y. (2019). Marginal Treatment Effects from a Propensity Score Perspective. \emph{JPE}, 127(6), 3070--3108. [PAP-zhou2019marginal]
\item Zhou, X., \& Xie, Y. (2020). Heterogeneous Treatment Effects in the Presence of Self-Selection. \emph{Sociological Methodology}, 50(1), 350--385. [PAP-zhou2020heterogeneous]
\end{itemize}

\textbf{Adaptive Experiments and Policy Learning (Section~\ref{sec:iii-adaptive}):}
\begin{itemize}[nosep]
\item Kasy, M., \& Sautmann, A. (2021). Adaptive Treatment Assignment in Experiments for Policy Choice. \emph{Econometrica}, 89(1), 113--132. [PAP-kasy2021adaptive]
\item Cesa-Bianchi, N., Colomboni, R., \& Kasy, M. (2025). Adaptive Maximization of Social Welfare. \emph{Econometrica}, forthcoming. [PAP-kasy2025adaptive]
\item Chernozhukov, V., et al. (2018). Double/Debiased ML for Treatment and Structural Parameters. \emph{Econometrics Journal}, 21(1), C1--C68. [PAP-chernozhukov2018double]
\item Athey, S., \& Imbens, G. (2016). Recursive Partitioning for Heterogeneous Causal Effects. \emph{PNAS}, 113(27), 7353--7360. [PAP-athey2016recursive]
\item Wager, S., \& Athey, S. (2018). Estimation and Inference of HTE using Random Forests. \emph{JASA}, 113(523), 1228--1242. [PAP-wager2018causalforest]
\item Caria, A. S., et al. (2024). An Adaptive Targeted Field Experiment: Job Search Assistance for Refugees in Jordan. \emph{JEEA}, 22(2), 781--836. [PAP-caria2024refugees]
\item Becker, G. M., DeGroot, M. H., \& Marschak, J. (1964). Measuring Utility by a Single-Response Sequential Method. \emph{Behavioral Science}, 9(3), 226--232. [PAP-becker1964bdm]
\item Vickrey, W. (1961). Counterspeculation, Auctions, and Competitive Sealed Tenders. \emph{Journal of Finance}, 16(1), 8--37. [PAP-vickrey1961counterspeculation]
\end{itemize}

\textbf{Causes of Effects and Retrospective Attribution (Section~\ref{sec:iii-coe}):}
\begin{itemize}[nosep]
\item Pearl, J. (2015). Causes of Effects and Effects of Causes. \emph{Sociological Methods \& Research}, 44(1), 149--164. [PAP-pearl2015causes]
\item Dawid, A. P., \& Musio, M. (2022). Effects of Causes and Causes of Effects. \emph{Annual Review of Statistics}, 9, 261--287. [PAP-dawid2022effects]
\item Tian, J., \& Pearl, J. (2000). Probabilities of Causation: Bounds and Identification. \emph{Ann. Math. AI}, 28, 287--313. [PAP-tian2000probabilities]
\item Luo, S., et al. (2025). Causal Attribution Analysis for Continuous Outcomes. \emph{ICML 2025}. [PAP-luo2025attribution]
\item Greenland, S., \& Robins, J. M. (2000). Epidemiology, Justice, and the Probability of Causation. \emph{Jurimetrics}, 40(3), 321--340. [PAP-greenland2000epidemiology]
\item Hart, H. L. A., \& Honor{\'e}, T. (1985). Causation in the Law (2nd ed.). OUP. [PAP-hart1985causation]
\end{itemize}

\textbf{Sequential Causal Inference (Section~\ref{sec:iii-seq}):}
\begin{itemize}[nosep]
\item Syrgkanis, V., \& Zhan, R. (2025). Post Reinforcement Learning Inference. \emph{Operations Research}, forthcoming. [PAP-syrgkanis2025postrl]
\item Zhang, Z., et al. (2024). Personalized Policy Learning through Discrete Experimentation. \emph{Management Science}. [PAP-zhang2024personalized]
\item Robins, J. M., et al. (2000). Marginal Structural Models and Causal Inference in Epidemiology. \emph{Epidemiology}, 11(5), 550--560. [PAP-robins2000msm]
\item Murphy, S. A. (2003). Optimal Dynamic Treatment Regimes. \emph{JRSS-B}, 65(2), 331--355. [PAP-murphy2003optimal]
\item Oaxaca, R. (1973). Male--Female Wage Differentials in Urban Labor Markets. \emph{IER}, 14(3), 693--709. [PAP-oaxaca1973decomposition]
\item Athey, S., \& Wager, S. (2021). Policy Learning with Observational Data. \emph{Econometrica}, 89(1), 133--161. [PAP-athey2021policy]
\end{itemize}

\textbf{Framework References:}
\begin{itemize}[nosep]
\item \texttt{docs/methods/alpha-definition.md} --- $\alpha$ conceptual definition
\item \texttt{docs/methods/alpha-vs-utility-formal.md} --- $\alpha$ vs. $U$ distinction
\item \texttt{docs/examples/green-energy-default-proof-of-concept.md} --- Full worked example
\end{itemize}

\textbf{Master Bibliography Reference:}

This appendix draws on literature indexed in:

\texttt{bibliography/bcm2\_master\_references.bib}

\nocite{bcm2_master}

% -----------------------------------------------------------------------------
% CROSS-REFERENCE FOOTER
% -----------------------------------------------------------------------------

\begin{tcolorbox}[colback=gray!5!white, colframe=gray!50!black]
\small
\textbf{Cross-References:}
Appendix HHH (METHOD-TOOLKIT) $\leftarrow$ uses R-Score |
Appendix DSN (METHOD-DESIGN) $\leftarrow$ uses calibrated $\alpha$ |
Appendix EST (CORE-WHERE) $\rightarrow$ provides anchors |
Appendix XX (METHOD-EVIDENCE) $\rightarrow$ evidence tier requirements for LLMMC input |
Appendix GLS (Glossary) $\rightarrow$ symbol definitions |
\textbf{Appendix PM (REF-PARAMETH)} $\leftarrow$ methodology navigation hub |
\textbf{Appendix DP (REF-DEVPIPE)} $\leftarrow$ Stage 3 (Calibration) in 6-stage pipeline |
\textbf{Appendix DR (REF-DATAREQ)} $\rightarrow$ data source inventory |
\textbf{Appendix VC (METHOD-VALIDATE)} $\leftarrow$ validates estimated parameters (DCV-3, DCV-4)
\end{tcolorbox}

% =============================================================================
% END OF APPENDIX III
% =============================================================================
